\scanpage{49}
$I$ ens ordonné

On va \unsure{interpréter|interprêter} $I^{\wedge}$ comme un $\operatorname{Top}(X)$

On a $I\simeq\operatorname{Top}(X)$ ; considérons la topologie canonique, Pour ceci, quels sont les cribles de $I$ : parties $\unsure{\sigma|\mathcal O|O}$ de $I$ telles que pour $\forall\ x\in\sigma$, les $y\leq x$ sont dans $\sigma$. On met sur $I$ la topologie définie par cet ens. de \deleted{cribles} parties (\deleted{\unsure{anti-|anti}}\unsure{appeler|appelle|pellicules} la topologie des cribles) On trouve que \deleted{\unsure{$\operatorname{Top}($|$\operatorname{Top}(I$}} l'on.
\[
\begin{aligned}
\operatorname{Ouv}(I)&\overset{\mathrm{déf}}{=}\operatorname{Crib}(I)\simeq\operatorname{Ouv}(I^{\wedge}) \qquad\text{donc}\\
{\unsure{I|\mathbb I}}^{\wedge}&\simeq\operatorname{Ouv}(I)^{\sim}=\operatorname{Top}(I).
\end{aligned}
\]

\uline{NB} La relation $y\leq x$ signifie que $y$ est une \uline{générisation} de $x$ (\unsure{non|ou|oui|$x$} : une spécialisation)

\diagram{\begin{tikzcd}[ampersand replacement=\&,row sep=large,nodes={anchor=west},column sep=0pt]
 \bullet\quad\parbox{.82\linewidth}{$E$ a suffisamment de \unsure{pts essentiels|pts essentiels,} et est engendré par ses ouverts} \arrow[d,leftrightarrow,start anchor=south west,end anchor=north west]\\
 \parbox{.86\linewidth}{$E$ est engendré par les \deleted{\unsure{objets|objet}} sous-objets de l'objet final qui sont connexes, non vides et projectifs.} \arrow[d,Leftrightarrow,start anchor=south west,end anchor=north west]\\
 \parbox{.86\linewidth}{(\deleted{\unsure{où|Que}}\unsure{si|où|Ex.} $E=C^{\wedge}$) Tous les morphismes de $C$ sont des monomorphismes} \arrow[d,Leftrightarrow,start anchor=south west,end anchor=north west]\\
 \parbox{.86\linewidth}{(si $E=\operatorname{Top}(X)$, $X$ sobre) Pour $\forall\ x\in X$, l'ens des générisations de $x$ est un voisinage de $x$.}
\end{tikzcd}}
% archive sheet number: 48
